\documentclass[12pt,a4paper]{article}

% -------------------------------------------------
% Pakete
% -------------------------------------------------
\usepackage{fontspec}
\usepackage[ngerman]{babel}
\usepackage{geometry}
\usepackage{setspace}
\usepackage{enumitem}
\usepackage{xcolor}
\usepackage{hyperref}
\usepackage{titlesec}
\usepackage{array}
\usepackage{booktabs}

% -------------------------------------------------
% Seitenlayout
% -------------------------------------------------
\geometry{
	left=2.5cm,
	right=2.5cm,
	top=2.5cm,
	bottom=2.5cm
}

\onehalfspacing
\setlength{\parindent}{0pt}
\setlength{\parskip}{0.6em}

% -------------------------------------------------
% Hyperlinks
% -------------------------------------------------
\hypersetup{
	colorlinks=true,
	linkcolor=black,
	urlcolor=blue,
	pdftitle={Kurzfassung Soll-Konzept},
	pdfauthor={},
	pdfsubject={Digitale Rückbestätigung geplanter Heizöl-Lieferungen}
}

% -------------------------------------------------
% Überschriften
% -------------------------------------------------
\titleformat{\section}
{\normalfont\Large\bfseries}
{\thesection}{1em}{}

\titleformat{\subsection}
{\normalfont\large\bfseries}
{\thesubsection}{1em}{}

% -------------------------------------------------
% Listen
% -------------------------------------------------
\setlist[itemize]{
	noitemsep,
	topsep=0.3em
}

% -------------------------------------------------
% Dokument
% -------------------------------------------------
\begin{document}
	
	% -------------------------------------------------
	% Titelseite
	% -------------------------------------------------
	\begin{titlepage}
		\centering
		\vspace*{3cm}
		
		{\Huge\bfseries Kurzfassung Soll-Konzept\par}
		\vspace{1cm}
		
		{\LARGE Digitale Rückbestätigung geplanter Heizöl-Lieferungen\par}
		\vspace{2cm}
		
		{\large Dokument zur fachlichen Abstimmung\par}
		
		\vfill
		
		{\large \today\par}
	\end{titlepage}
	
	\tableofcontents
	\newpage
	
	% -------------------------------------------------
	\section{Ziel der Lösung}
	
	Ziel der geplanten Lösung ist es, den manuellen Aufwand bei der Rückbestätigung geplanter Heizöl-Lieferungen zu reduzieren.
	
	Die bestehende Tourplanung bleibt dabei unverändert. Das System soll die Disposition nicht ersetzen, sondern den Disponenten bei der strukturierten Einholung, Erfassung und Auswertung von Kundenrückmeldungen unterstützen.
	
	Die Kernaussage lautet:
	
	\begin{quote}
		Die Lösung digitalisiert die Rückbestätigung geplanter Lieferfenster, ohne die Disposition selbst zu automatisieren.
	\end{quote}
	
	% -------------------------------------------------
	\section{Ausgangssituation}
	
	Der Disponent plant eine Tour wie bisher im DISPO-System. Dabei werden Aufträge einer Tour zugeordnet und Lieferfenster festgelegt.
	
	Anschließend muss geprüft werden, ob die vorgeschlagenen Liefertermine für die Kunden passen. Dieser Schritt erfolgt aktuell überwiegend manuell, zum Beispiel telefonisch.
	
	Dadurch entstehen mehrere Probleme:
	
	\begin{itemize}
		\item viele manuelle Anrufversuche;
		\item Kunden sind telefonisch schwer erreichbar;
		\item Rückrufe verzögern die Planung;
		\item mehrere Mitarbeiter können in die Klärung eingebunden sein;
		\item der aktuelle Rückmeldestatus ist nicht immer sofort übersichtlich.
	\end{itemize}
	
	% -------------------------------------------------
	\section{Soll-Prozess}
	
	Im Soll-Prozess werden bestehende Aufträge aus dem DISPO-System übernommen und dem Disponenten angezeigt.
	
	Das Rückbestätigungssystem erzeugt keine eigenen Aufträge und verändert die Tourplanung nicht automatisch.
	
	Nachdem eine Tour vorläufig geplant wurde, kann der Disponent für einzelne Aufträge eine Rückbestätigungsanfrage auslösen.
	
	Die Anfrage wird also nicht automatisch und nicht pauschal für die gesamte Tour versendet, sondern bewusst pro Auftrag.
	
	Eine Rückbestätigungsanfrage bezieht sich immer auf:
	
	\begin{itemize}
		\item einen konkreten Auftrag;
		\item ein konkretes Lieferfenster;
		\item eine konkrete Anfrage.
	\end{itemize}
	
	Der Kunde erhält eine Nachricht mit den für ihn relevanten Informationen:
	
	\begin{itemize}
		\item Lieferdatum;
		\item Lieferfenster;
		\item Lieferadresse;
		\item Bestellmenge.
	\end{itemize}
	
	Der Kunde kann das vorgeschlagene Lieferfenster bestätigen oder ablehnen. Zusätzlich kann er optional einen Kommentar hinterlassen.
	
	Dieser Kommentar wird nicht automatisch interpretiert und führt nicht automatisch zu einer neuen Planung. Er dient nur als Hinweis für den Disponenten.
	
	% -------------------------------------------------
	\section{Rückbestätigungsstatus}
	
	Für einen Auftrag sind folgende Rückbestätigungsstatus vorgesehen:
	
	\begin{itemize}
		\item \textbf{offen}
		\item \textbf{versendet}
		\item \textbf{bestätigt}
		\item \textbf{abgelehnt}
		\item \textbf{keine Rückmeldung}
	\end{itemize}
	
	\subsection{Offen}
	
	\textbf{Offen} bedeutet, dass für den Auftrag noch keine aktuelle Rückbestätigungsanfrage erfolgreich versendet wurde.
	
	\subsection{Versendet}
	
	\textbf{Versendet} bedeutet, dass die Anfrage technisch erfolgreich verschickt wurde und das System auf eine Kundenantwort wartet.
	
	\subsection{Bestätigt}
	
	\textbf{Bestätigt} bedeutet, dass der Kunde das Lieferfenster innerhalb der gültigen Frist bestätigt hat.
	
	\subsection{Abgelehnt}
	
	\textbf{Abgelehnt} bedeutet, dass der Kunde das Lieferfenster innerhalb der gültigen Frist abgelehnt hat.
	
	\subsection{Keine Rückmeldung}
	
	\textbf{Keine Rückmeldung} bedeutet, dass innerhalb von 24 Stunden keine gültige Antwort eingegangen ist.
	
	% -------------------------------------------------
	\section{Wichtige fachliche Regeln}
	
	Für die Rückbestätigung gelten folgende fachliche Regeln:
	
	\begin{itemize}
		\item Der Status \textbf{versendet} wird erst gesetzt, wenn der technische Versand erfolgreich war.
		\item Ein Klick auf \glqq Senden\grqq{} reicht dafür nicht aus.
		\item Die Antwortfrist beträgt 24 Stunden.
		\item Die Frist startet mit erfolgreichem Versand der Anfrage.
		\item Pro Anfrage ist nur eine gültige Antwort möglich.
		\item Statuswirksam ist immer nur die aktuellste Anfrage.
		\item Antworten auf alte, überholte oder abgelaufene Anfragen ändern den aktuellen Status nicht automatisch.
		\item Wenn das Lieferfenster nach dem Versand geändert wird, muss für das neue Lieferfenster eine neue Anfrage erzeugt werden.
		\item Bei Ablehnung oder fehlender Rückmeldung erstellt das System keine automatische Neuplanung.
		\item Das System schlägt keine automatischen Ersatztermine vor.
	\end{itemize}
	
	% -------------------------------------------------
	\section{Problemfälle und manuelle Bearbeitung}
	
	Aufträge mit dem Status \textbf{abgelehnt} oder \textbf{keine Rückmeldung} werden dem Disponenten besonders angezeigt.
	
	Diese Aufträge erscheinen zum Beispiel in einer Arbeitsliste:
	
	\begin{quote}
		Manuell zu bearbeiten
	\end{quote}
	
	Der Disponent entscheidet dann selbst, wie weiter verfahren wird.
	
	Mögliche Optionen sind:
	
	\begin{itemize}
		\item neues Lieferfenster festlegen und neue Anfrage senden;
		\item Kunden telefonisch kontaktieren;
		\item Auftrag auf eine andere Tour verschieben;
		\item Fall anderweitig dispositiv entscheiden.
	\end{itemize}
	
	Wenn ein Fall telefonisch geklärt wurde, kann der Disponent den Auftrag manuell auf \textbf{bestätigt} setzen.
	
	Ein zusätzlicher Status wie \textbf{manuell bestätigt} ist nicht vorgesehen, da dieser keine eigene Kundenantwort beschreibt, sondern nur die Art der Klärung.
	
	% -------------------------------------------------
	\section{Tourstatus und Freigabe}
	
	Die Rückbestätigung erfolgt auf Auftragsebene.
	
	Die Tour dient als Übersicht und Gruppierung mehrerer Aufträge.
	
	Der Status \textbf{final geplant} gehört nicht zur Rückbestätigung, sondern zur Tourplanung.
	
	Ein Auftrag kann also bestätigt sein, ohne dass die gesamte Tour final geplant ist.
	
	Für Touren können separate Planungsstatus verwendet werden, zum Beispiel:
	
	\begin{itemize}
		\item \textbf{vorläufig geplant}
		\item \textbf{freigabebereit}
		\item \textbf{final geplant}
	\end{itemize}
	
	Eine Tour kann als \textbf{freigabebereit} angezeigt werden, wenn alle relevanten Aufträge geklärt sind.
	
	Geklärt bedeutet zum Beispiel:
	
	\begin{itemize}
		\item der Auftrag wurde digital bestätigt;
		\item der Auftrag wurde telefonisch durch den Disponenten geklärt;
		\item der Auftrag wurde auf eine andere Tour verschoben;
		\item der Fall wurde anderweitig dispositiv entschieden.
	\end{itemize}
	
	Die finale Freigabe der Tour bleibt eine bewusste Entscheidung des Disponenten.
	
	% -------------------------------------------------
	\section{Geplante Oberfläche}
	
	Die Oberfläche soll dem Disponenten eine zentrale Übersicht bieten.
	
	Vorgesehen sind:
	
	\begin{itemize}
		\item eine Auftragssicht mit allen relevanten Aufträgen;
		\item eine Tourübersicht mit allen Aufträgen einer Tour;
		\item eine Arbeitsliste für problematische Fälle;
		\item Detailinformationen zu einzelnen Aufträgen;
		\item Aktionen wie \glqq Anfrage senden\grqq, \glqq Neue Anfrage senden\grqq, \glqq Telefonisch geklärt\grqq{} oder \glqq Auftrag verschieben\grqq.
	\end{itemize}
	
	Die Oberfläche soll nicht einzelne E-Mails in den Mittelpunkt stellen, sondern den fachlichen Status der Aufträge.


	
\end{document}